\documentclass[../../../hsm_tok_q.tex]{subfiles}

\begin{document}
    \setcounter{figure}{0}
    Kさんのクラスでは，先生が示した問題をみんなで考えた．

    次の各問に答えよ．

    \begin{mycolorboxa}[title=［先生が示した問題］]
        \textsf{図\ref{fig/hsm_tok_2025_2nd_02_01_q}}に示した立体は，%
        底面が半径2\:cmの円，高さが5\:cmの円柱で，底面の2つの円の中心をそれぞれO，Pとし，%
        点Oと点Pを結んでできる線分は2つの底面に垂直である．
       
        円Oの周上にある点をQとし，点Oと点P，点Oと点Qをそれぞれ結び，点Qを通り線分OPに平行な直線を引き，%
        円Pの周上で交わる点をRとし，点Pと点Rを結ぶ．
        
        四角形OPRQの面積を$\mathrm{S\:cm^2}$，円柱の体積を$\mathrm{V\:cm^3}$とするとき，%
        VをSを用いた式で表しなさい．
        \\
        \begin{minipage}{\linewidth}
            \vspace{.7\baselineskip}
            {\centering
                \setlength{\abovecaptionskip}{5pt}
                \includegraphics%
                    {\subfix{fig_hsm_tok_2025_2nd_02/fig_hsm_tok_2025_2nd_02_01_q.pdf}}%
                    \captionof{figure}{} \label{fig/hsm_tok_2025_2nd_02_01_q}%
            }%
        \end{minipage}
    \end{mycolorboxa}
    %
    Kさんは，［先生が示した問題］の答えを次の形の式で表した．%
    Kさんの答えは正しかった．

    \vspace{.2\baselineskip}
    〈Kさんの答え〉 $\mathrm{V}\jequal\:\myboxf{　}\:\mathrm{S}$
        
    \begin{enumerate}
        \item 〈Kさんの答え〉の\:\myboxf{　}\:に当てはまるものを，%
            次の{\textsf ア}〜\textsf{エ}のうちから選び，記号で答えよ．

            ただし，円周率は$\pi$とする．

            {\renewcommand{\arraystretch}{1.5}%
                \vspace{-.2\baselineskip}
                \begin{tabularx}{\dimexpr\linewidth-\zw}{@{}LLLL@{}}
                    {\textsf ア}\; $2\pi$ & {\textsf イ}\; $4\pi$%
                    & {\textsf ウ}\; $10\pi$ & {\textsf エ}\; $20\pi$
                \end{tabularx}
            }
    \end{enumerate}

    \newpage
    
    Kさんのグループは，［先生が示した問題］をもとにして，次の問題を考えた．

    \begin{mycolorboxa}[title=［Kさんのグループが作った問題］]
        $a$，$h$を正の数とする．

        \textsf{図\ref{fig/hsm_tok_2025_2nd_02_02_q}}に示した立体ABCD--EFGHは，%
        底面が1辺$2a\:\mathrm{cm}$の正方形，高さが$h\:\mathrm{cm}$の直方体である．

        四角形ABCDにおいて，対角線ACと対角線BDとの交点をO，四角形EFGHにおいて，
        対角線EGと対角線FHとの交点をPとする．

        辺BC，辺FGの中点をそれぞれQ，Rとし，点Oと点P，点Oと点Q，点Pと点R，点Qと点Rをそれぞれ結ぶ．

        四角形ABCDの周の長さを$\ell\:\mathrm{cm}$，四角形OPRQの面積を$\mathrm{S\:cm^2}$，%
        立体ABCD--EFGHの体積を$\mathrm{V\:cm^3}$とするとき，%
        $\mathrm{V}=\myfraca{1}{2}\ell\mathrm{S}$となることを確かめてみよう．
        \\
        \begin{minipage}{\linewidth}
            \vspace{.5\baselineskip}
            {\centering
            \setlength{\abovecaptionskip}{2pt}
                \includegraphics%
                    {\subfix{fig_hsm_tok_2025_2nd_02/fig_hsm_tok_2025_2nd_02_02_q.pdf}}%
                    \captionof{figure}{} \label{fig/hsm_tok_2025_2nd_02_02_q}%
            }%
        \end{minipage}
        %
    \end{mycolorboxa}
    
    \begin{enumerate}[resume]
        \item ［Kさんのグループが作った問題］で，%
            $\mathrm{V}=\myfraca{1}{2}\ell\mathrm{S}$となることを証明せよ．
    \end{enumerate}
\end{document}
